% registry: FIN
\methodhead{FIN}{Findley (1959): shear amplitude plus maximum normal stress on the critical plane}\label{sheet:FIN}

\block{Source} W.~N.~Findley, \emph{J.\ Eng.\ Ind.} (1959) \cite{findley1959}, Eqs.~(2), (8)--(10),
(13), (18), (19), (24), pp.~302--305.

\block{Idea and scope} Fatigue damage on a material plane is driven by the alternating shear stress
on that plane, increased linearly by the largest normal stress acting on it. The critical plane is
the plane where this combination is largest. Proposed for high-cycle fatigue of metals under
combined bending (or axial load) and torsion with mean stresses; the paper treats proportional
(in-phase) loading. The repository evaluates constant-amplitude, single-harmonic loading.

\block{Inputs}
\begin{center}\small
\begin{tabular}{@{}llp{100mm}@{}}
\toprule
Symbol & Unit & Meaning\\\midrule
$\theta$ & rad & orientation of the material plane (normal measured from the greatest principal stress in Eq.~(13))\\
$\tau_{\theta A}=\tau_a$ & MPa & shear stress amplitude on the plane\\
$\sigma_{\theta M}=\sigma_{n,\max}$ & MPa & maximum normal stress on the plane over the cycle, $\sigma_{n,m}+\sigma_{n,a}$ (Eq.~\eqref{eq:scalar})\\
$\sigma_{n,m}$, $\sigma_{n,a}$ & MPa & mean and amplitude of the normal stress on the plane\\
$\Delta\tau$ & MPa & shear stress range on the plane, $\tau_a=\Delta\tau/2$\\
$a\ge b$ & MPa & semi-axes of the elliptical shear path of an out-of-phase cycle\\
$k$ ($k_\mathrm{FI}$) & -- & normal stress sensitivity, material constant (workbook \texttt{k FI})\\
$f$ & MPa & Findley fatigue strength at the life considered\\
$\sigma_A$, $\tau_A$ & MPa & bending and torsion stress amplitudes (in-phase loading)\\
$\sigma_M$, $\tau_M$ & MPa & bending and torsion maximum stresses (in-phase loading)\\
$\fm$, $\tm$ & MPa & fully reversed bending and torsion fatigue limits\\
$r=\tm/\fm$ & -- & limit ratio, $\tfrac12<r<1$\\
$\mathrm{DP}$ & MPa & repository damage parameter, Eq.~\eqref{eq:fin-dp}\\
\bottomrule
\end{tabular}
\end{center}

\block{Formulation} On the plane of orientation $\theta$ (Eqs.~2, 8--10, p.~302),
\begin{equation}
\tau_{\theta A} + k\,\sigma_{\theta M} = f ,
\label{eq:fin-main}
\end{equation}
and the critical plane is the one that maximises the left-hand side. For in-phase bending and
torsion the maximum over the planes is (Eq.~18)
\begin{equation}
f = \sqrt{\Bigl(\frac{\sigma_A}{2}\Bigr)^2 + \tau_A^2 + k^2\Bigl[\Bigl(\frac{\sigma_M}{2}\Bigr)^2 + \tau_M^2\Bigr]} + k\,\frac{\sigma_M}{2} .
\label{eq:fin-eq18}
\end{equation}
Fully reversed torsion (Eq.~19) and fully reversed bending (Eq.~24) give
\begin{equation}
f = \tm\sqrt{1+k^2},\qquad f = \frac{\fm}{2}\left(k+\sqrt{1+k^2}\right),
\end{equation}
so the limit ratio fixes $k$:
\begin{equation}
\frac{\tm}{\fm} = \frac{k+\sqrt{1+k^2}}{2\sqrt{1+k^2}}
\quad\Longleftrightarrow\quad
k = \frac{2r-1}{\sqrt{1-(2r-1)^2}} ,\qquad r=\frac{\tm}{\fm}\in\left(\tfrac12,1\right).
\label{eq:fin-k}
\end{equation}
$k\to0$ for $r\to\tfrac12$ and $k\to\infty$ for $r\to1$. With zero mean stress the critical plane lies
at $\tan2\theta=1/k$ from the greatest principal stress (Eq.~13). In the repository, on every
candidate plane,
\begin{equation}
\mathrm{DP} = \tau_a + k_\mathrm{FI}\,\sigma_{n,\max},\qquad \tau_a=\tfrac12\Delta\tau,\qquad
\sigma_{n,\max}=\sigma_{n,m}+\sigma_{n,a} .
\label{eq:fin-dp}
\end{equation}
\emph{Critical plane}: the plane of largest DP (the selection metric is DP itself, so \texttt{FIN} and
\texttt{FIN\_ext} coincide); ties go to the larger damage (Section~\ref{sec:planes}).
\emph{Shear amplitude}: proportional loading gives $\tau_a=a$; out of phase, the default rule is
\texttt{rss}, $\tau_a=\sqrt{a^2+b^2}$ (MCE). The global selection of \texttt{docs/shear-amplitude.md}
can switch it to the longest chord (MCC, $\tau_a=a$).

\block{Constants and calibration} Paper: $k$ from the ratio of the fully reversed torsion and bending
limits, Eq.~\eqref{eq:fin-k}, or from fatigue strengths in two loading modes at the same life.
Repository: $k_\mathrm{FI}$ from the workbook, a \texttt{-{}-constants} record, or the pipeline
calibration. By default the pipeline searches $k_\mathrm{FI}\in[0,3]$ for the value that minimises
the standard deviation of $\ln N_f$ of the $R=-1$ tension and torsion failures about one master
curve, re-selecting the critical plane for each trial $k$; these tests become design rows. On
fatigue-limit datasets $k_\mathrm{FI}$ makes the $R=-1$ tension and torsion limits give the same DP,
which is the original procedure.

\block{Life} The 1959 paper prescribes no life law. The repository uses the master curve on DP
(Section~\ref{sec:mastercurve}); $\mathrm{DP}/f$ is a fatigue index.

\block{Repository choices}
\begin{itemize}
\item $\sigma_{n,\max}$ keeps its sign: a compressive normal stress lowers DP; no clipping.
\item Out-of-phase shear amplitude: \texttt{rss} of the shear-ellipse semi-axes; the paper gives no
  rule. A $k_\mathrm{FI}$ calibrated with the MCC amplitude is rejected by the production (MCE)
  assessment.
\item Life from the shared master curve $\mathrm{DP}=AN_f^b$, not a Findley-specific law.
\item $k_\mathrm{FI}$ calibrated by default on the life scatter of tension and torsion failures, not
  on the fatigue-limit ratio; it reduces to Eq.~\eqref{eq:fin-k} only if the two S--N curves are
  parallel and the limit ratio holds at every life.
\item The mean shear stress does not enter DP, as in the original criterion.
\end{itemize}

\block{Worked example} Synthetic material, $\fm=300$, $\tm=200\MPa$: Eq.~\eqref{eq:fin-k} gives
$k_\mathrm{FI}=0.3536$ and $f=\tm\sqrt{1+k^2}=212.13\MPa$. The plane search returns
$\mathrm{DP}=212.13\MPa$ at both limits. Bending with $\sigma_a=200$, $\sigma_m=100\MPa$
($\sigma_M=300\MPa$) gives $\mathrm{DP}=166.23\MPa$, as Eq.~\eqref{eq:fin-eq18}, a fatigue index
of 0.784; without the mean the DP would be $100\,(k+\sqrt{1+k^2})=141.4\MPa$.

\block{Verification} The plane search reproduces the closed forms Eqs.~(18), (19), (24) and the
critical-plane angle of Eq.~(13) for $k=0.1$, 0.2, 0.3, 1.1, and the limit ratios of p.~303. The
angles printed on p.~305 differ from Eq.~(13) by up to $0.65^\circ$ and between loadings; they appear
read from a figure (paper inconsistency). Also tested: analytical plane values, MCC/MCE amplitudes,
calibration and the explicit-amplitude path. Not reproduced: the experimental data of the paper or
an independent multiaxial dataset.
